\documentclass[12pt]{article}
\usepackage[utf8]{inputenc}
\usepackage[T1]{fontenc}
\usepackage{amsmath,amssymb,amsthm}
\usepackage{graphicx}
\usepackage[margin=1in]{geometry}
\usepackage{hyperref}
\usepackage{booktabs}
\usepackage{natbib}
\usepackage[nopatch=footnote]{microtype}

\newtheorem{theorem}{Theorem}
\newtheorem{lemma}[theorem]{Lemma}
\newtheorem{corollary}[theorem]{Corollary}
\newtheorem{proposition}[theorem]{Proposition}
\theoremstyle{definition}
\newtheorem{definition}[theorem]{Definition}
\theoremstyle{remark}
\newtheorem{remark}[theorem]{Remark}

\newcommand{\Hil}{\mathcal{H}}
\newcommand{\Alg}{\mathcal{A}}
\newcommand{\Tr}{\operatorname{Tr}}
\newcommand{\Inn}{\operatorname{Inn}}
\newcommand{\Aut}{\operatorname{Aut}}
\newcommand{\Out}{\operatorname{Out}}
\newcommand{\ad}{\operatorname{ad}}
\newcommand{\Lie}{\mathcal{L}}
\newcommand{\KK}{M_{\mathrm{KK}}}

\title{Ungaugeability of Kosmann Symmetries\\of the Dirac Operator on Compact Lie Groups}
\author{MemeTheory\thanks{AI assistance in research and manuscript preparation (Anthropic Claude: Opus 4.6--4.8, Fable 5).}}
\date{\today}

\begin{document}

\maketitle

\begin{abstract}
We prove that the $U(1)_7$ symmetry generated by the Kosmann derivative $K_7$ of the Dirac operator on Jensen-deformed $\mathrm{SU}(3)$ cannot be gauged within Connes' noncommutative geometry framework.  Working with the almost-commutative spectral triple $\Alg = C^\infty(\mathrm{SU}(3)) \otimes \Alg_F$ with $\Alg_F = \mathbb{C} \oplus \mathbb{H} \oplus M_3(\mathbb{C})$ (the Standard Model algebra), we show that $K_7$ is excluded from the gauge sector by two structurally distinct arguments.  First, the commutant obstruction: $[K_7, D_K(\tau)] = 0$ for all Jensen parameter values $\tau$, so $K_7$ generates a trivial one-form in $\Omega^1_D(\Alg)$.  By the spectral theorem, this extends to all functionals of $D_K$, closing the obstruction non-perturbatively.  Second, a categorical argument independent of the commutant property: $K_7$ generates a diffeomorphism of $\mathrm{SU}(3)$ and therefore belongs to $\Aut(C^\infty(\mathrm{SU}(3)))$; since $C^\infty(\mathrm{SU}(3))$ is commutative, $\Inn(C^\infty(\mathrm{SU}(3))) = \{\mathrm{id}\}$ and all its automorphisms are outer; gauge fields arise from $\Inn(\Alg_F)$, not from diffeomorphisms of the base.  The sharp form of the statement is that inner automorphisms of $\Alg_F$ -- which generate the Standard Model gauge group -- \emph{can} be gauged, while diffeomorphisms of $C^\infty(\mathrm{SU}(3))$, including the Kosmann symmetry $K_7$, \emph{cannot}; $U(1)_7$ lies on the diffeomorphism side of this divide.  The inner/outer distinction is classical in Connes' framework; what is new is an explicit worked example demonstrating ungaugeability of a Kosmann symmetry in a concrete Kaluza--Klein setting, with computational verification, together with the observation that the same commutant identity is simultaneously the signature of an isometry and the obstruction to its gauging.  The Anderson--Higgs impossibility for $U(1)_7$ follows as a direct corollary; we contrast it with the complementary case of non-Killing internal fields, whose gauge fields acquire mass without an ad hoc Higgs.
\end{abstract}

%=====================================================================
\section{Introduction}
\label{sec:intro}
%=====================================================================

A central achievement of Connes' noncommutative geometry (NCG) is the derivation of gauge theories from purely geometric data \cite{Connes1996gravity, ChamseddineConnes1996spectral, vanSuijlekom2015book}.  Given a real spectral triple $(\Alg, \Hil, D, J, \gamma)$, gauge fields arise as \emph{inner fluctuations} of the Dirac operator:
\begin{equation}\label{eq:innerfluc}
  D \;\longmapsto\; D_A = D + A + \varepsilon' J A J^{-1},
  \qquad
  A = \sum_i a_i [D, b_i],
  \quad a_i, b_i \in \Alg.
\end{equation}
The self-adjoint one-form $A$ belongs to the space $\Omega^1_D(\Alg)$ of \emph{universal one-forms} evaluated through the representation of $D$.  The gauge group is the quotient $\mathcal{G} = U(\Alg)/U(Z(\Alg))$, where $U(\Alg)$ denotes the unitary group of $\Alg$ and $Z(\Alg)$ its center \cite{Connes1996gravity, vanSuijlekom2015book}.  In the almost-commutative geometry $M^4 \times F$ that recovers the Standard Model, the gauge group emerges as $U(1)_Y \times \mathrm{SU}(2)_L \times \mathrm{SU}(3)_c$, and the Higgs field appears as the component of the inner fluctuation along the finite direction $F$ \cite{CCM2007gravity}.  More precisely, when the inner fluctuation \eqref{eq:innerfluc} is decomposed by the action of the relevant isometry group, the one-form splits automatically into a spin-1 part along the \emph{Killing} directions -- the gauge fields -- and a spin-0 part along the \emph{non-Killing} directions -- the Higgs.  This Killing/non-Killing dichotomy is central to what follows: the symmetry $U(1)_7$ we study is generated by a Killing field, so it can only ever supply a spin-1 (gauge) one-form, never a Higgs; and the commutant property will show that even that spin-1 one-form is trivial.

This framework draws a sharp line between two classes of symmetries.  \emph{Inner automorphisms} of $\Alg$ -- conjugation by unitaries $u \in U(\Alg)$ -- generate gauge transformations and produce nontrivial gauge fields through the map $b \mapsto [D, b]$.  \emph{Outer automorphisms} -- elements of $\Aut(\Alg)/\Inn(\Alg)$ -- correspond to diffeomorphisms of the underlying space and do \emph{not} generate gauge fields.  This inner/outer dichotomy, established in Connes' foundational work \cite{Connes1994book, Connes1996gravity}, is the noncommutative-geometric realization of the physical distinction between gauge symmetries and spacetime symmetries.

The question we address is: \emph{can a symmetry of the Dirac operator that arises from a diffeomorphism of the underlying manifold ever be promoted to a gauge symmetry through any mechanism available within the spectral action framework?}

We answer this question in the negative for a specific, concrete case.  Consider the Dirac operator $D_K$ on $\mathrm{SU}(3)$ equipped with a one-parameter family of left-invariant metrics -- the Jensen deformation \cite{Jensen1973}, parametrized by $\tau \geq 0$.  At $\tau = 0$, the metric is bi-invariant and the full $\mathrm{SU}(3) \times \mathrm{SU}(3)$ isometry group acts on the spinor bundle through Kosmann derivatives \cite{Kosmann1971}.  For $\tau > 0$, the Jensen deformation breaks the right-multiplication isometries down to the stabilizer of the deformation direction; the residual right-isometry contains the subgroup $U(1)_7$ generated by the seventh Gell-Mann matrix $\lambda_7$ (see Section~\ref{subsec:jensen} for the precise stabilizer).  The corresponding Kosmann derivative $K_7$ continues to commute with $D_K(\tau)$ at every $\tau$:
\begin{equation}\label{eq:commutant}
  [iK_7, \, D_K(\tau)] = 0, \qquad \forall\, \tau \geq 0.
\end{equation}
On the singlet sector this identity holds exactly (the commutator is the zero operator to floating-point precision); the analytic proof, valid for all $\tau$, is given in Proposition~\ref{prop:commutant}.

Despite being an exact symmetry of the Dirac operator, $U(1)_7$ cannot be gauged.  We establish this through two structurally distinct arguments, each sufficient on its own:

\begin{enumerate}
  \item \textbf{Commutant obstruction}: $K_7$ lies in the commutant of $D_K$, hence $a[D_K, K_7] = 0$ for all $a$ in any algebra $\Alg$.  The one-form $\Omega^1_D(\Alg)$ in the $K_7$ direction is trivial.  By the spectral theorem, this extends non-perturbatively to all functionals of $D_K$, including quantum corrections at every loop order.

  \item \textbf{Categorical}: $K_7$ generates a diffeomorphism of $\mathrm{SU}(3)$, which is an outer automorphism of $C^\infty(\mathrm{SU}(3))$ (since all automorphisms of a commutative algebra are outer).  In the almost-commutative spectral triple $\Alg = C^\infty(\mathrm{SU}(3)) \otimes \Alg_F$, gauge fields arise from $\Inn(\Alg_F)$, not from diffeomorphisms of the base.  This argument does \emph{not} use $[K_7, D_K] = 0$.
\end{enumerate}

The physical consequence is that $U(1)_7$ admits no Anderson--Higgs mechanism: there is no gauge boson to absorb the Goldstone mode of $U(1)_7$ breaking, no mass generation through the standard Higgs mechanism, and no gauge kinetic term $F_7^{\mu\nu} F_{7,\mu\nu}$ in the spectral action.

The distinction between inner and outer automorphisms in NCG is well established \cite{Connes1994book, Connes1996gravity, ConnesMarcolli2008book}.  For the commutative algebra $C^\infty(\mathrm{SU}(3))$ alone, ungaugeability of diffeomorphisms follows from general principles.  What is new in this work is threefold: an explicit worked example demonstrating ungaugeability of a Kosmann symmetry in a concrete Kaluza--Klein setting -- the Jensen-deformed $\mathrm{SU}(3)$ with an almost-commutative spectral triple -- verified computationally and established by two structurally independent arguments; the observation that the all-orders extension of the commutant obstruction (Corollary~\ref{cor:allorders}) is a single statement subsuming what is naively split into tree-level, one-loop, and non-perturbative obstructions; and the recognition that the commutant identity \eqref{eq:commutant} plays a double role -- it is at once the Kosmann--Lichnerowicz signature of an isometry and the trivial-one-form obstruction to gauging that isometry.


%=====================================================================
\section{Preliminaries}
\label{sec:prelim}
%=====================================================================

\subsection{Real spectral triples}
\label{subsec:spectraltriple}

We recall the essential definitions following Connes \cite{Connes1995reality, Connes1996gravity}.

\begin{definition}[Real spectral triple]
A \emph{real spectral triple} of KO-dimension $n \bmod 8$ is a quintuple $(\Alg, \Hil, D, J, \gamma)$ where:
\begin{itemize}
  \item $\Alg$ is a unital $*$-algebra represented faithfully on a separable Hilbert space $\Hil$;
  \item $D$ is a self-adjoint operator on $\Hil$ with compact resolvent, such that $[D, a]$ is bounded for all $a \in \Alg$;
  \item $J : \Hil \to \Hil$ is an antiunitary operator (the \emph{real structure});
  \item $\gamma$ is a self-adjoint unitary ($\mathbb{Z}/2$-grading) when $n$ is even;
\end{itemize}
subject to the sign conditions
\begin{equation}\label{eq:KOsigns}
  J^2 = \varepsilon, \quad JD = \varepsilon' DJ, \quad J\gamma = \varepsilon'' \gamma J,
\end{equation}
where $(\varepsilon, \varepsilon', \varepsilon'') \in \{+1, -1\}^3$ are determined by $n \bmod 8$ according to the KO-dimension table \cite{Connes1995reality}.
\end{definition}

The KO-dimension relevant to the Standard Model is $n = 6 \bmod 8$, corresponding to $(\varepsilon, \varepsilon', \varepsilon'') = (+1, +1, -1)$.

\subsection{The almost-commutative spectral triple on \texorpdfstring{$\mathrm{SU}(3)$}{SU(3)}}
\label{subsec:almostcommutative}

The paper works throughout with an \emph{almost-commutative} spectral triple, which we now specify.  The purely commutative algebra $C^\infty(\mathrm{SU}(3))$ has a trivial inner automorphism group: for any commutative algebra $\Alg$, every automorphism $\alpha(a) = u a u^* = a$ is the identity when $\Alg$ is commutative, so $\Inn(\Alg) = \{\mathrm{id}\}$.  Consequently, the gauge group of a purely commutative spectral triple is trivial and no symmetry can be gauged -- making any ungaugeability result vacuous.  To obtain a nontrivial gauge sector, we work with the product geometry:

\begin{definition}[Almost-commutative spectral triple on $\mathrm{SU}(3)$]\label{def:almostcomm}
\sloppy The spectral triple is $(\Alg, \Hil, D, J, \gamma)$ where:
\begin{itemize}
  \item The algebra is $\Alg = C^\infty(\mathrm{SU}(3)) \otimes \Alg_F$ with finite algebra $\Alg_F = \mathbb{C} \oplus \mathbb{H} \oplus M_3(\mathbb{C})$;
  \item The Hilbert space is $\Hil = L^2(\mathrm{SU}(3), \Sigma) \otimes \Hil_F$ where $\Hil_F = \mathbb{C}^{32}$ is the Standard Model fermion space (one generation, particle and antiparticle);
  \item The Dirac operator is $D = D_K(\tau) \otimes \mathbf{1}_F + \gamma_9 \otimes D_F$, where $D_K(\tau)$ is the Dirac operator on $(\mathrm{SU}(3), g_\tau)$, $\gamma_9$ is the chirality operator on $\mathrm{SU}(3)$, and $D_F$ is the finite Dirac operator encoding fermion masses and mixings;
  \item $J = J_{\mathrm{SU}(3)} \otimes J_F$ and $\gamma = \gamma_9 \otimes \gamma_F$ are the product real structure and grading.
\end{itemize}
\end{definition}

The gauge group of this spectral triple is
\begin{equation}\label{eq:gaugegroupSM}
  \mathcal{G} = \Inn(\Alg) = \Inn(\Alg_F) \cong U(1)_Y \times \mathrm{SU}(2)_L \times \mathrm{SU}(3)_c,
\end{equation}
where the equality $\Inn(\Alg) = \Inn(\Alg_F)$ follows from the commutativity of $C^\infty(\mathrm{SU}(3))$: inner automorphisms can only arise from the noncommutative factor $\Alg_F$.  The Standard Model gauge fields -- the hypercharge, weak, and color gauge bosons -- emerge as inner fluctuations of $D$ from this algebra; concretely, the unimodular unitaries of $\Alg_F$ furnish exactly the thirteen generators of $U(1) \times \mathrm{SU}(2) \times \mathrm{SU}(3)$.

The Kosmann derivative $K_7$ acts on the \emph{first} tensor factor: it is an automorphism of $C^\infty(\mathrm{SU}(3))$, not of $\Alg_F$.  Since $C^\infty(\mathrm{SU}(3))$ is commutative, $K_7 \in \Aut(C^\infty(\mathrm{SU}(3)))$ is outer -- in fact, \emph{every} nontrivial automorphism of a commutative algebra is outer.  Even in the product $\Alg = C^\infty(\mathrm{SU}(3)) \otimes \Alg_F$, the element $K_7 \otimes \mathbf{1}_F$ acts trivially on $\Alg_F$ and as an outer automorphism on $C^\infty(\mathrm{SU}(3))$, so it remains outer for the full algebra.  This makes the ungaugeability result nontrivial: the SM gauge symmetries CAN be gauged (they come from $\Inn(\Alg_F)$), but $U(1)_7$ CANNOT (it comes from a diffeomorphism of $\mathrm{SU}(3)$).

\begin{remark}[Two distinct symmetry routes]\label{rem:tworoutes}
It is worth separating two groups that both involve $\mathrm{SU}(3)$ but play structurally different roles.  The gauge group \eqref{eq:gaugegroupSM} is the \emph{unitary} group of the finite algebra $\Alg_F$; it sources the inner fluctuations.  The isometry group of the Jensen metric $g_\tau$ on the \emph{manifold} $\mathrm{SU}(3)$ is a \emph{different} group -- it supplies the Peter--Weyl labels of the spectrum and the Kosmann derivatives, but it is not the chiral gauge group, and the two are not interchangeable.  The symmetry $U(1)_7$ studied here belongs to the second (isometry) route, which is precisely why it cannot be gauged.
\end{remark}

\subsection{Inner automorphisms and gauge fields}
\label{subsec:innerauto}

The automorphism group $\Aut(\Alg)$ of the algebra acts on the spectral triple by transforming $D \mapsto \alpha(D)$ for $\alpha \in \Aut(\Alg)$.  The subgroup of \emph{inner automorphisms} $\Inn(\Alg)$ consists of conjugations $a \mapsto u a u^*$ by unitaries $u \in U(\Alg)$.

In Connes' framework \cite{Connes1996gravity, Connes1994book, vanSuijlekom2015book}, the gauge group of the spectral triple is
\begin{equation}\label{eq:gaugegroup}
  \mathcal{G} = \Inn(\Alg) \,\cong\, U(\Alg) / U(Z(\Alg)),
\end{equation}
where $Z(\Alg)$ denotes the center of $\Alg$.  Gauge connections are inner fluctuations of the Dirac operator \eqref{eq:innerfluc}, and the space of gauge potentials is the space of self-adjoint one-forms
\begin{equation}\label{eq:omega1}
  \Omega^1_D(\Alg) = \Big\{\sum_i a_i [D, b_i] : a_i, b_i \in \Alg\Big\}.
\end{equation}

The fundamental principle is that an element $b \in \Alg$ generates a nontrivial gauge field if and only if $[D, b] \neq 0$.  Elements in the commutant $\{b \in \Alg : [D, b] = 0\}$ generate \emph{zero} gauge fields.

\begin{remark}[Inner vs.\ outer]
Outer automorphisms -- elements of $\Out(\Alg) = \Aut(\Alg)/\Inn(\Alg)$ -- correspond to diffeomorphisms of the underlying space in the commutative case ($\Alg = C^\infty(M)$, where $\Aut(\Alg) \cong \mathrm{Diff}(M)$).  These do \emph{not} generate gauge fields.  The full automorphism group of a spectral triple decomposes as
\[
  \Aut(\Alg) = \Inn(\Alg) \rtimes \Out(\Alg) \;\longleftrightarrow\; \text{gauge} \times \text{diffeos}.
\]
This is the noncommutative-geometric formulation of the equivalence principle: gravity (diffeomorphisms) and gauge forces (inner automorphisms) are geometrically unified but categorically distinct.
\end{remark}

\subsection{The spectral action}
\label{subsec:spectralaction}

The bosonic action functional is the spectral action \cite{ChamseddineConnes1996spectral}:
\begin{equation}\label{eq:spectralaction}
  S_b = \Tr\, f\!\left(\frac{D_A^2}{\Lambda^2}\right),
\end{equation}
where $f : \mathbb{R}_+ \to \mathbb{R}_+$ is a smooth, positive, decreasing cutoff function and $\Lambda$ is an energy scale.  The asymptotic expansion for $\Lambda \to \infty$ yields
\begin{equation}\label{eq:SAexpansion}
  S_b \sim 2 f_4 \Lambda^4 a_0 + 2 f_2 \Lambda^2 a_2 + f_0 a_4 + O(\Lambda^{-2}),
\end{equation}
where $f_{2k} = \int_0^\infty f(u)\, u^{k-1}\, du$ are the moments of $f$, and $a_{2k}$ are the Seeley--DeWitt coefficients of $D_A^2$.  For the almost-commutative Standard Model geometry, $a_0$ gives the cosmological constant, $a_2$ gives the Einstein--Hilbert action with a Higgs mass term, and $a_4$ gives the Yang--Mills, Higgs quartic, and gravitational curvature-squared terms \cite{CCM2007gravity}.

A gauge kinetic term $\frac{1}{g_a^2} \int |F_a|^2\, d\mathrm{vol}$ for a gauge field $A_a$ appears in $a_4$ if and only if the corresponding generator produces a nontrivial one-form through \eqref{eq:omega1}.  If $[D, b] = 0$ for the generator $b$, no such term arises.

\subsection{The Kosmann derivative}
\label{subsec:kosmann}

Let $(M, g)$ be a Riemannian manifold with spin structure, and let $X$ be a Killing vector field.  The \emph{Kosmann derivative} \cite{Kosmann1971} is the canonical lift of $X$ to an operator on the spinor bundle $\Sigma M$:
\begin{equation}\label{eq:kosmann}
  K_X = \nabla_X^{\Sigma} + \tfrac{1}{4}\, (\nabla_a X_b)\, \gamma^a \gamma^b,
\end{equation}
where $\nabla^{\Sigma}$ is the spin connection and $\gamma^a$ are the Clifford generators in an orthonormal frame.  When $X$ is a Killing field, the Killing equation $\nabla_a X_b + \nabla_b X_a = 0$ guarantees that the second term is antisymmetric in $a, b$ and thus involves only the $\gamma^{[a}\gamma^{b]}$ component of the Clifford product.  The Kosmann derivative satisfies the Leibniz rule with respect to Clifford multiplication and preserves the spin structure.  For a Killing field, $K_X$ commutes with the Dirac operator:
\begin{equation}\label{eq:kosmanncommutes}
  [K_X, D] = 0 \qquad \text{whenever } X \text{ is Killing}.
\end{equation}
This is a classical result; the spinor Lie derivative is due to Kosmann \cite{Kosmann1971}, building on the harmonic-spinor and transformation-group analysis of Lichnerowicz \cite{Lichnerowicz1963}, and the commutation property is treated in the modern spin-geometry framework of Fatibene and Francaviglia \cite{Fatibene1996}.

The Kosmann derivative generates an infinitesimal isometry of the spin Riemannian manifold.  It is a first-order differential operator on sections of $\Sigma M$, encoding the action of the isometry group on spinors.  In the language of NCG, the Kosmann derivative is a derivation of the algebra $C^\infty(M)$ that preserves the Dirac operator -- but it is an \emph{outer} derivation, not an inner one, because diffeomorphisms of $M$ are outer automorphisms of $C^\infty(M)$.

\subsection{Jensen deformation of \texorpdfstring{$\mathrm{SU}(3)$}{SU(3)}}
\label{subsec:jensen}

The Jensen deformation \cite{Jensen1973} is a one-parameter family of left-invariant Riemannian metrics on a compact Lie group $G$, defined as follows.  Fix a closed subgroup $H \subset G$ and let $\mathfrak{g} = \mathfrak{h} \oplus \mathfrak{m}$ be the reductive decomposition of the Lie algebra.  The Jensen metric $g_\tau$ is defined on $\mathfrak{g}$ by
\begin{equation}\label{eq:jensenmetric}
  g_\tau(X, Y) =
  \begin{cases}
    e^{-4\tau}\, B(X, Y) & \text{if } X, Y \in \mathfrak{h}, \\
    e^{\tau}\, B(X, Y) & \text{if } X, Y \in \mathfrak{m}, \\
    0 & \text{if } X \in \mathfrak{h},\; Y \in \mathfrak{m},
  \end{cases}
\end{equation}
where $B$ is the Killing form (normalized to be positive definite) and $\tau \geq 0$ is the deformation parameter.  At $\tau = 0$, the metric is proportional to the bi-invariant Killing metric.  For $\tau > 0$, the metric remains left-invariant but breaks right-invariance.  The deformation is volume-preserving: $\det(g_\tau) / \det(g_0) = 1$ for all $\tau$, the exponential rescalings of $\mathfrak{h}$ and $\mathfrak{m}$ being balanced so that the total Riemannian volume of $\mathrm{SU}(3)$ is fixed (to machine precision in the explicit computation).

For $G = \mathrm{SU}(3)$ and $H = U(1)$ generated by the seventh Gell-Mann matrix $\lambda_7$, the deformation breaks the bi-invariant isometry group.  Writing $\mathrm{Iso}(g_0) = (\mathrm{SU}(3)_L \times \mathrm{SU}(3)_R)/\mathbb{Z}_3$ for the bi-invariant metric, the left factor is preserved for any left-invariant metric, while the right factor is broken to the stabilizer of the deformation direction:
\begin{equation}\label{eq:isometrybreaking}
  \frac{\mathrm{SU}(3)_L \times \mathrm{SU}(3)_R}{\mathbb{Z}_3}
  \;\xrightarrow{\;\tau > 0\;}\;
  \mathrm{SU}(3)_L \times \big(\mathrm{SU}(2)_R \times U(1)_7\big),
\end{equation}
where $\mathrm{SU}(2)_R \times U(1)_7 \cong U(2)_R$ is the centralizer of $\lambda_7$ in $\mathrm{SU}(3)_R$ (a twelve-dimensional residual right-isometry).  This residual stabilizer is the pattern identified by Baptista \cite{Baptista2024KK} for the unravelling of the bi-invariant $\mathrm{SU}(3)$ metric, in which the isometry group descends to $(\mathrm{SU}(3) \times \mathrm{SU}(2) \times U(1))/\mathbb{Z}_6$.  The subgroup of \eqref{eq:isometrybreaking} relevant to the present paper is the abelian factor $U(1)_7 \subset U(2)_R$ generated by $\lambda_7$; it is on this $U(1)_7$ that the Kosmann derivative $K_7$ acts, and whose ungaugeability we establish.


%=====================================================================
\section{The \texorpdfstring{$U(1)_7$}{U(1)\_7} Symmetry of \texorpdfstring{$D_K$}{D\_K}}
\label{sec:U1symmetry}
%=====================================================================

Let $D_K(\tau)$ denote the Dirac operator on $(\mathrm{SU}(3), g_\tau)$ for the Jensen metric \eqref{eq:jensenmetric}.  The Kosmann derivative $K_7$ associated to the Killing field $\lambda_7$ satisfies \eqref{eq:kosmanncommutes} at $\tau = 0$ by the general theory.  The nontrivial content is that this commutation relation survives the Jensen deformation.

\begin{proposition}\label{prop:commutant}
For all $\tau \geq 0$,
\begin{equation}\label{eq:commutantexact}
  [iK_7, \, D_K(\tau)] = 0.
\end{equation}
\end{proposition}

\begin{proof}
The Killing field $X_7 = \lambda_7$ generates a one-parameter subgroup that remains an isometry of $g_\tau$ for all $\tau$, since $\lambda_7$ generates the subgroup $H = U(1)$ that defines the Jensen deformation \eqref{eq:jensenmetric}.  The decomposition $\mathfrak{su}(3) = \mathfrak{h} \oplus \mathfrak{m}$ is $\ad(\lambda_7)$-invariant by construction.  The Jensen metric $g_\tau$ rescales $\mathfrak{h}$ and $\mathfrak{m}$ by different factors, but this rescaling commutes with $\ad(\lambda_7)$ because $\lambda_7 \in \mathfrak{h}$ and $\ad(\lambda_7)$ preserves each summand.  Therefore $X_7$ remains Killing for $g_\tau$ at every $\tau$, and the general identity \eqref{eq:kosmanncommutes} applies, giving $[iK_7, D_K(\tau)] = 0$ for all $\tau \geq 0$.
\end{proof}

\begin{remark}[Numerical illustration]\label{rem:numillustration}
The identity \eqref{eq:commutantexact} was confirmed by direct computation on the sixteen-dimensional singlet sector of the Peter--Weyl decomposition of $L^2(\mathrm{SU}(3), \Sigma)$, where $iK_7$ and $D_K(\tau)$ are explicit $16 \times 16$ matrices.  At each of nine deformation values $\tau \in [0, 0.50]$ the commutator $[iK_7, D_K(\tau)]$ is the zero matrix to floating-point precision: its Frobenius norm vanishes identically (the relative norm $\|[iK_7, D_K(\tau)]\|_F / \|D_K(\tau)\|_F$ is $0$ to the working precision of the computation).  The same machinery, applied to the full-spectrum compatibility identity $[J, D_K(\tau)] = 0$ across all $79{,}968$ matrix-element pairs of the larger representation, returns a maximal relative deviation below $3.3 \times 10^{-13}$, corroborating that the commutation structure of $D_K(\tau)$ with the relevant symmetry operators is exact at the level of the algebra.
\end{remark}

\begin{remark}[Charge decomposition]
The operator $iK_7$ has eigenvalues $q \in \{0, \pm\tfrac{1}{4}\}$ on the singlet sector.  This induces a charge decomposition of the Dirac eigenspace into three branches:
\begin{itemize}
  \item $B_1$: charge $q = 0$ (four modes, $U(2)$ singlet),
  \item $B_2$: charge $q = \pm\tfrac{1}{4}$ (four modes, $U(2)$ doublet),
  \item $B_3$: charge $q = 0$ (eight modes, $U(2)$ adjoint).
\end{itemize}
The four-, four-, and eight-mode branches exhaust the sixteen-dimensional singlet sector.  The charge eigenvalues remain exactly $\pm 1/4$ to machine precision at all $\tau$.  All spectral selection rules -- including the vanishing of the pairing kernel $V(B_1, B_1) = 0$ -- follow from $U(1)_7$ charge conservation on these branches.
\end{remark}


%=====================================================================
\section{Two Proofs of Ungaugeability}
\label{sec:proofs}
%=====================================================================

We establish the main result through two structurally distinct arguments.  The first exploits the commutant property $[K_7, D_K] = 0$; the second is independent of it.

\begin{theorem}[Ungaugeability of $U(1)_7$]\label{thm:main}
Let $(\Alg, \Hil, D, J, \gamma)$ be the almost-commutative spectral triple of Definition~\ref{def:almostcomm} on $\mathrm{SU}(3)$ with Jensen metric $g_\tau$.  Then the $U(1)_7$ symmetry generated by $K_7$ cannot be gauged: there exists no inner fluctuation $A \in \Omega^1_D(\Alg)$ that produces a nontrivial gauge field for $U(1)_7$.  Consequently, $U(1)_7$ admits no Anderson--Higgs mechanism within this spectral triple.
\end{theorem}

\subsection{Proof I: Commutant obstruction}
\label{subsec:proof1}

\begin{proof}[Proof I]
By Proposition~\ref{prop:commutant}, $K_7$ lies in the commutant of $D_K(\tau)$:
\[
  K_7 \in \{T \in \mathcal{B}(\Hil) : [D_K, T] = 0\}.
\]
The inner fluctuation in the $K_7$ direction, for any $a \in \Alg$, is
\begin{equation}\label{eq:trivialoneform}
  A_7 = a \cdot [D_K, K_7] = a \cdot 0 = 0.
\end{equation}
Therefore the component of $\Omega^1_D(\Alg)$ along $K_7$ is trivial: $K_7$ generates the zero one-form.

The spectral action $\Tr f(D_A^2/\Lambda^2)$ expanded to second order in a gauge potential $A$ produces the gauge kinetic term through the Seeley--DeWitt coefficient $a_4$.  This coefficient contains terms of the form $\Tr(F_{\mu\nu}^a F^{a\,\mu\nu})$, where $F^a_{\mu\nu}$ is the field strength of the gauge field $A^a_\mu = \sum_i a_i [D, b_i]_\mu$ associated to the generator $b$ of the gauge symmetry.  Since $[D_K, K_7] = 0$, we have $A_7 = 0$, $F_7 = 0$, and no kinetic term appears.  The $U(1)_7$ gauge boson is absent from the spectral action.
\end{proof}

\begin{corollary}[Non-perturbative extension]\label{cor:allorders}
The commutant obstruction holds at all orders in perturbation theory and extends non-perturbatively.
\end{corollary}

\begin{proof}
If $[K_7, D_K] = 0$, then $[K_7, p(D_K)] = 0$ for every polynomial $p$, by induction:
\begin{equation*}
[K_7, D_K^{n+1}] = [K_7, D_K^n]\, D_K + D_K^n\, [K_7, D_K] = 0.
\end{equation*}
More generally, $[K_7, f(D_K)] = 0$ for every measurable function $f$: since $K_7$ commutes with $D_K$, it commutes with each spectral projection $E_\lambda$ of the spectral measure $\{E_\lambda\}$ (by the double commutant theorem applied to the von Neumann algebra generated by $D_K$), and hence with $f(D_K) = \int f(\lambda)\, dE_\lambda$.

The one-loop counterterms of the spectral action have the form $\Sigma(D_K) = c_0 + c_2 D_K^2 + c_4 D_K^4 + \cdots$ \cite{vanNulandvanSuijlekom2022}, so the effective Dirac operator $D_{\mathrm{eff}} = D_K + \Sigma(D_K)$ satisfies
\begin{equation}\label{eq:oneloopcommutant}
  [iK_7, D_{\mathrm{eff}}] = [iK_7, D_K] + [iK_7, \Sigma(D_K)] = 0 + 0 = 0.
\end{equation}
The same argument applies at every loop order, and the spectral theorem extends it to arbitrary (non-polynomial) functionals of $D_K$.  The one-form obstruction $A_7 = a[D_{\mathrm{eff}}, K_7] = 0$ therefore holds non-perturbatively.  This single statement -- that $K_7$ commutes with every functional of $D_K$ -- subsumes what would otherwise be three separate obstructions (tree-level, one-loop, all-orders): each is a corollary of the one commutant identity \eqref{eq:commutantexact}.

The non-perturbative propagation can be checked explicitly on the singlet sector.  Writing $\Lambda = 3\, \KK$ for the cutoff used in $e^{-D_K^2/\Lambda^2}$, the relative Frobenius norms
\begin{align}
  \|[iK_7, D_K^2]\|_F / \|D_K^2\|_F &= 1.3 \times 10^{-17}, \notag \\
  \|[iK_7, D_K^4]\|_F / \|D_K^4\|_F &= 3.0 \times 10^{-17}, \label{eq:numverify} \\
  \|[iK_7, e^{-D_K^2/\Lambda^2}]\|_F / \|e^{-D_K^2/\Lambda^2}\|_F &= 3.1 \times 10^{-16}, \notag
\end{align}
are all at machine epsilon, confirming that the commutant property carries to the relevant functionals of $D_K$ at the level of the explicit operators.
\end{proof}

\subsection{Proof II: Categorical (inner vs.\ outer)}
\label{subsec:proof2}

This proof establishes the impossibility from the structure theory of automorphisms, \emph{without} using the commutant property $[K_7, D_K] = 0$.  It applies regardless of whether $K_7$ commutes with the Dirac operator.

\begin{proof}[Proof II]
The argument has four steps.

\medskip\noindent\textbf{Step 1: $K_7$ generates a diffeomorphism.}
The Kosmann derivative $K_7$ is the canonical lift \eqref{eq:kosmann} of the Killing vector field $X_7 = \lambda_7$ to the spinor bundle $\Sigma(\mathrm{SU}(3))$.  The flow $\phi_t = \exp(tX_7)$ is a one-parameter group of isometries of $(\mathrm{SU}(3), g_\tau)$, and $K_7$ generates its lift to the spin structure.  In particular, $K_7$ generates a one-parameter group of diffeomorphisms acting on sections of $\Sigma(\mathrm{SU}(3))$.

\medskip\noindent\textbf{Step 2: Diffeomorphisms are automorphisms of $C^\infty(\mathrm{SU}(3))$.}
A diffeomorphism $\phi : \mathrm{SU}(3) \to \mathrm{SU}(3)$ induces an algebra automorphism $\alpha_\phi : C^\infty(\mathrm{SU}(3)) \to C^\infty(\mathrm{SU}(3))$ by pullback: $\alpha_\phi(f) = f \circ \phi$.  The Gelfand--Naimark theorem identifies $\Aut(C^\infty(M)) \cong \mathrm{Diff}(M)$ for compact manifolds $M$ \cite{Connes1994book}.

\medskip\noindent\textbf{Step 3: For commutative algebras, all automorphisms are outer.}
For a commutative unital $*$-algebra $\Alg$, conjugation by any unitary $u \in U(\Alg)$ is trivial: $u a u^* = u u^* a = a$ for all $a \in \Alg$.  Therefore $\Inn(\Alg) = \{\mathrm{id}\}$, and every nontrivial automorphism is outer:
\begin{equation}\label{eq:inncommutative}
  \Alg \text{ commutative} \quad\Longrightarrow\quad \Out(\Alg) = \Aut(\Alg), \quad \Inn(\Alg) = \{\mathrm{id}\}.
\end{equation}
In particular, the automorphism $\alpha_{\phi_t}$ induced by the flow of $K_7$ is outer for $C^\infty(\mathrm{SU}(3))$.

\medskip\noindent\textbf{Step 4: Gauge fields come from $\Inn(\Alg_F)$, not from diffeomorphisms.}
In the almost-commutative spectral triple of Definition~\ref{def:almostcomm}, $\Alg = C^\infty(\mathrm{SU}(3)) \otimes \Alg_F$.  The gauge group is $\mathcal{G} = \Inn(\Alg) = \Inn(\Alg_F) \cong U(1)_Y \times \mathrm{SU}(2)_L \times \mathrm{SU}(3)_c$.  Gauge connections are inner fluctuations $A = \sum_i a_i [D, b_i]$ with $a_i, b_i \in \Alg$, and the gauge fields that participate in the Anderson--Higgs mechanism are precisely those arising from $\Inn(\Alg_F)$.

The generator $K_7 \otimes \mathbf{1}_F$ acts as a diffeomorphism of $\mathrm{SU}(3)$ tensored with the identity on $\Alg_F$.  It is an outer automorphism of the first factor and the identity on the second.  It is therefore outer for the product algebra $\Alg$.  Since gauge fields arise exclusively from inner automorphisms \cite{Connes1996gravity}, $K_7$ cannot produce a gauge field.
\end{proof}

\begin{remark}[Independence of the two proofs]
Proof~I uses $[K_7, D_K] = 0$ and is independent of the algebra structure beyond the existence of $\Omega^1_D(\Alg)$.  Proof~II uses only the categorical position of $K_7$ as a diffeomorphism generator and the commutativity of $C^\infty(\mathrm{SU}(3))$ -- it does \emph{not} invoke $[K_7, D_K] = 0$ at any step.  Proof~II would hold even for a hypothetical Killing field whose Kosmann derivative did \emph{not} commute with $D_K$.  The two arguments are therefore structurally distinct.  The numerical computations of Remark~\ref{rem:numillustration} and equation \eqref{eq:numverify} are not a third, independent proof: they corroborate Proof~I, whose content is already established analytically.
\end{remark}

\begin{corollary}[Anderson--Higgs impossibility]\label{cor:noanderson}
The $U(1)_7$ symmetry of $D_K(\tau)$ on Jensen-deformed $\mathrm{SU}(3)$ admits no Anderson--Higgs mechanism within the almost-commutative spectral triple of Definition~\ref{def:almostcomm}.  Specifically:
\begin{enumerate}
  \item There exists no gauge boson $A_7^\mu$ in the spectral action for $U(1)_7$.
  \item The Goldstone mode of $U(1)_7$ breaking cannot be absorbed.
  \item No mass is generated for the $U(1)_7$ sector through the Higgs mechanism.
\end{enumerate}
\end{corollary}

\begin{remark}[The impossibility is specific to the Killing--commutant sector]\label{rem:notgeneral}
Corollary~\ref{cor:noanderson} is not a statement that mass generation is impossible in the spectral triple of Definition~\ref{def:almostcomm}; it is specific to $U(1)_7$, a symmetry generated by a Killing field that lies in the commutant of $D_K$.  Recall from Section~\ref{sec:intro} that the inner fluctuation \eqref{eq:innerfluc} decomposes into a spin-1 part along the Killing directions (gauge fields) and a spin-0 part along the non-Killing directions (the Higgs).  The Higgs field of the spectral action therefore lives in the \emph{non-Killing} sector; it is structurally disjoint from $U(1)_7$, which is Killing.  Thus the algebra admits a Higgs mechanism in general -- its absence for $U(1)_7$ is precisely because $U(1)_7$ sits in the wrong (Killing, commutant) sector to couple to it.  We return to this complementarity in Section~\ref{subsec:nonkilling}.
\end{remark}

\begin{remark}[$D_{\mathrm{phys}}$ and the fluctuated commutant]\label{rem:dphys}
The ungaugeability theorem (Theorem~\ref{thm:main}) applies to the \emph{unfluctuated} Dirac operator $D_K$.  When inner fluctuations from $\Alg_F$ are turned on, the fluctuated operator
\[
  D_{\mathrm{phys}} = D_K + A + \varepsilon' J A J^{-1}, \qquad A \in \Omega^1_D(\Alg),
\]
may have a different commutant.  Concretely, at gap-edge fluctuation amplitude $\phi = 0.133$ (in $\KK$ units), direct computation on the singlet sector gives
\[
  \frac{\|[iK_7, D_{\mathrm{phys}}]\|_F}{\|D_{\mathrm{phys}}\|_F} = 0.117 .
\]
This $11.7\%$ breaking of $[iK_7, D_{\mathrm{phys}}] = 0$ shows that $K_7$ does \emph{not} lie in the commutant of $D_{\mathrm{phys}}$, and one might ask whether this opens a route to gauging $U(1)_7$.

It does not.  The breaking arises from inner fluctuations of $\Alg_F$ -- gauge fields for $\mathrm{SU}(2)_L \times U(1)_Y$, not for $U(1)_7$.  Gaugeability in NCG requires $K_7 \in \Inn(\Alg)$, which is a categorical property of $K_7$'s position in the automorphism group (Proof~II).  Whether $K_7$ commutes with $D_{\mathrm{phys}}$ or not is irrelevant: $K_7$ generates a diffeomorphism of $\mathrm{SU}(3)$, which is outer for $C^\infty(\mathrm{SU}(3))$, and no fluctuation of $D$ can promote an outer automorphism to an inner one.  The commutant breaking by $D_{\mathrm{phys}}$ means that $U(1)_7$ ceases to be an \emph{exact symmetry} of the fluctuated operator, but it was never a \emph{gauge symmetry} to begin with.
\end{remark}


%=====================================================================
\section{Consequences}
\label{sec:consequences}
%=====================================================================

\subsection{The Goldstone boson persists}
\label{subsec:goldstone}

In condensed matter and particle physics, when a continuous symmetry is spontaneously broken, Goldstone's theorem guarantees the existence of a massless boson.  The Anderson--Higgs mechanism \cite{Anderson1963, Higgs1964} provides the sole known escape: if the broken symmetry is gauged, the Goldstone boson is absorbed as the longitudinal polarization of the massive gauge boson.

For $U(1)_7$ on Jensen-deformed $\mathrm{SU}(3)$, a BCS-type condensate formed from eigenmodes of $D_K$ carries $K_7$ charge $\pm 1/2$ and spontaneously breaks $U(1)_7$.  The Goldstone mode associated to this breaking is realized as a sharp collective (Leggett-type) mode of the condensate, with frequency $\omega_G \approx 0.070\, \KK$.  By Theorem~\ref{thm:main}, no gauge boson exists to absorb it.  The Goldstone mode is permanent and ungapped within the NCG inner-fluctuation framework.

\subsection{Implications for Kaluza--Klein spectral triples}
\label{subsec:kkimplications}

In a Kaluza--Klein geometry $M^4 \times K$ with internal compact manifold $K$, the isometry group $\mathrm{Iso}(K)$ acts on the spectrum of the internal Dirac operator $D_K$ through Kosmann derivatives.\footnote{We use the standard Kaluza--Klein language of ``the internal manifold $K$'' throughout; in a spectral-geometric reading the internal data are intrinsic to the triple $(C^\infty(K), L^2(K,\Sigma), D_K)$ -- the spectrum of $D_K$ \emph{is} the internal content -- rather than fields propagating on a pre-existing container.  The two readings agree on every statement made here.}  The Jensen deformation provides a smooth family of metrics on $K = \mathrm{SU}(3)$ that breaks $\mathrm{Iso}(K)$ from $(\mathrm{SU}(3) \times \mathrm{SU}(3))/\mathbb{Z}_3$ to $\mathrm{SU}(3)_L \times U(2)_R$ as in \eqref{eq:isometrybreaking}.

Theorem~\ref{thm:main} shows that the residual abelian isometry $U(1)_7 \subset U(2)_R$ is \emph{structurally different} from the gauge symmetries $\mathrm{SU}(3)_c \times \mathrm{SU}(2)_L \times U(1)_Y$ that arise from the algebra $\Alg_F$.  The isometry symmetry and the gauge symmetry belong to categorically distinct classes of automorphisms -- outer and inner, respectively -- and no mechanism within the spectral action framework can convert one into the other.

This has a broader implication, which we record as the structural core of the paper.

\begin{proposition}[Isometry--obstruction duality]\label{prop:duality}
Let $K$ be a compact spin manifold and $X$ a Killing vector field for a metric $g$ on $K$, with Kosmann lift $K_X$ on $\Sigma K$.  In any Kaluza--Klein spectral triple whose internal Dirac operator is $D_K$ for $(K, g)$, the identity $[K_X, D_K] = 0$ is simultaneously
\begin{enumerate}
  \item the Kosmann--Lichnerowicz signature that $X$ is an isometry, and
  \item the obstruction to gauging $X$: it forces $A_X = a[D_K, K_X] = 0$, so $K_X$ generates the trivial one-form in $\Omega^1_{D_K}$.
\end{enumerate}
Consequently, any internal isometry whose Kosmann derivative commutes with $D_K$ is automatically ungaugeable.
\end{proposition}

\begin{proof}
Statement (1) is \eqref{eq:kosmanncommutes}.  Statement (2) is \eqref{eq:trivialoneform} applied to $K_X$.  The two are the same equation $[K_X, D_K] = 0$ read in two registers: as a property of the flow of $X$ (isometry) and as a property of the one-form map $a \mapsto a[D_K, K_X]$ (gauging).  The final clause is the conjunction of (1) and (2).
\end{proof}

For $X = \lambda_7$ on $\mathrm{SU}(3)$ this is Theorem~\ref{thm:main}; Proposition~\ref{prop:duality} isolates the part of the argument that holds for any compact group manifold and any commutant Killing field.

\subsection{Non-Killing fields: the complementary case}
\label{subsec:nonkilling}

The ungaugeability established here is the Killing-side half of a dichotomy.  Baptista \cite{Baptista2024KK, Baptista2025chiral} has analyzed the complementary case of gauge fields associated to \emph{non-Killing} fields on the internal manifold $K$ -- fields that preserve the Einstein--Hilbert action on $K$ but not the full metric.  Such fields are spontaneously broken by the choice of vacuum metric and generate \emph{massive} gauge bosons, with mass controlled by the Lie derivative of the vacuum metric, and with no need for an ad hoc Higgs field; the bosons can be arbitrarily light and can couple chirally to fermions, evading the standard no-go arguments against chiral fermionic interactions in Kaluza--Klein theory \cite{Baptista2024KK, Baptista2025chiral}.

The two results bracket the same structural statement about $\mathrm{Iso}(K)$ versus $\mathrm{Diff}(K)$:
\begin{itemize}
  \item \emph{Killing field, Kosmann derivative in the commutant of $D_K$} (this paper): $[K_X, D_K] = 0$, the one-form is trivial, the symmetry is ungaugeable, and the associated Goldstone mode is not absorbed.
  \item \emph{Non-Killing field} (Baptista): $\Lie_X g \neq 0$, the gauge boson is massive without an external Higgs, and chiral couplings are possible.
\end{itemize}
In the spectral-action reading of Remark~\ref{rem:notgeneral}, these are exactly the spin-1 (Killing) and spin-0/massive (non-Killing) sectors of the inner fluctuation.  The Higgs of the almost-commutative model lives on the non-Killing side; $U(1)_7$ lives on the Killing--commutant side, which is why it has neither a gauge boson nor a Higgs partner.

\subsection{Spectral action and the inner/outer dichotomy}
\label{subsec:saouter}

The spectral action principle \cite{ChamseddineConnes1996spectral} asserts that the entire bosonic Lagrangian is encoded in $\Tr f(D^2/\Lambda^2)$.  Theorem~\ref{thm:main} reveals a structural feature of this principle: the spectral action \emph{automatically} excludes gauge fields for symmetries generated by commutant elements.  This is not a defect of the formalism but a reflection of its internal consistency: the spectral action generates dynamics only for the degrees of freedom that appear as inner fluctuations, and the commutant directions are, by definition, the directions in which $D$ does not fluctuate.

The inner/outer dichotomy thus acquires a dynamical meaning through the spectral action: inner automorphisms produce propagating gauge fields with kinetic terms in $a_4$; outer automorphisms (including Kosmann symmetries in the commutant) produce global symmetries with conservation laws but no propagating bosons.


%=====================================================================
\section{Relation to Prior Work}
\label{sec:priorwork}
%=====================================================================

\subsection{Connes' inner/outer distinction}

The classification of automorphisms into inner (gauge) and outer (diffeomorphisms) is a foundational element of Connes' program, established in \cite{Connes1994book, Connes1996gravity} and reviewed in \cite{Connes2019spectral, ConnesMarcolli2008book, vanSuijlekom2015book}.  The gauge group derivation for the Standard Model -- showing that $\Inn(\Alg_F) \cong U(1)_Y \times \mathrm{SU}(2)_L \times \mathrm{SU}(3)_c$ for $\Alg_F = \mathbb{C} \oplus \mathbb{H} \oplus M_3(\mathbb{C})$ -- is one of the central results of NCG applied to particle physics \cite{CCM2007gravity}.  Our work does not extend or modify this framework; rather, we apply it to a specific manifold and symmetry to derive a concrete impossibility result.

\subsection{Kosmann's construction}

The Lie derivative of spinors was defined by Kosmann in 1971 \cite{Kosmann1971}, building on Lichnerowicz's transformation-group and harmonic-spinor analysis \cite{Lichnerowicz1963}, and was later cast in the modern spin-geometry / natural-bundle language by Fatibene and Francaviglia \cite{Fatibene1996}.  The key properties -- the Leibniz rule with respect to Clifford multiplication, commutation with the Dirac operator for Killing fields, and the structure \eqref{eq:kosmann} as a first-order differential operator on sections of the spinor bundle -- are classical.  Recent work by Baptista \cite{Baptista2024KK, Baptista2025chiral} has placed the Kosmann derivative in the context of Kaluza--Klein theory, showing how it governs the coupling of massive gauge fields to internal spinor modes; as discussed in Section~\ref{subsec:nonkilling}, his non-Killing analysis is the natural complement to the present Killing-field result.  We note an additional structural distinction: the Kosmann derivative is a \emph{first-order differential operator} (involving $\nabla_X^\Sigma$ and $\nabla_a X_b$), while inner fluctuations $a[D,b]$ are \emph{zeroth-order} (multiplication operators, since $[D, f] = -i\gamma^\mu \partial_\mu f$ is zeroth-order in derivatives when $f \in C^\infty(M)$).  This difference in operator order provides an independent perspective on why Kosmann derivatives cannot arise as inner fluctuations, complementing the categorical argument of Proof~II.

\subsection{Inner fluctuations without the first-order condition}

Chamseddine, Connes, and van Suijlekom \cite{CCS2013inner} extended the inner fluctuation framework to spectral triples that violate the first-order condition $[[D, a], b^\circ] = 0$.  In this generalized setting, quadratic terms $\sum c_{ij} [D, a_i][D, a_j]$ appear in the fluctuated Dirac operator.  The Dirac operator $D_K(\tau)$ on Jensen-deformed $\mathrm{SU}(3)$ violates the first-order condition, attaining the maximal Clifford value $4.000$ for $(\mathbb{H}, \mathbb{H})$ pairs.  One might ask whether the quadratic fluctuations could circumvent the commutant obstruction.

They cannot: the quadratic term involves products $[D_K, a_i][D_K, a_j]$ with $a_i, a_j \in \Alg$.  Since $[D_K, K_7] = 0$, any product involving $K_7$ as one of the arguments vanishes.  The commutant obstruction persists even in the generalized fluctuation framework.

\subsection{What is new}
\label{subsec:whatisnew}

The inner/outer distinction is known.  The Kosmann derivative is known.  The spectral action framework is known.  For the commutative algebra $C^\infty(\mathrm{SU}(3))$, the ungaugeability of any diffeomorphism follows from general principles ($\Inn = \{\mathrm{id}\}$).  What is new in this paper is:

\begin{enumerate}
  \item An \emph{explicit worked example} demonstrating ungaugeability of a Kosmann symmetry in a concrete Kaluza--Klein setting -- the Jensen-deformed $\mathrm{SU}(3)$ with an almost-commutative spectral triple -- with computational verification of $[iK_7, D_K(\tau)] = 0$ across the full deformation family.

  \item The identification that the commutant property and the spectral theorem together close the obstruction \emph{non-perturbatively} (Corollary~\ref{cor:allorders}) in a single statement, rather than as three separate tree-level/one-loop/all-orders facts.

  \item The \emph{isometry--obstruction duality} (Proposition~\ref{prop:duality}): the identity $[K_X, D_K] = 0$ is at once the Kosmann--Lichnerowicz signature of an isometry and the trivial-one-form obstruction to gauging it, so that on any compact group manifold a commutant internal isometry is automatically ungaugeable.

  \item The \emph{physical consequence} and its boundary: the impossibility of the Anderson--Higgs mechanism for $U(1)_7$, the persistence of the Goldstone mode, and the demarcation (Remarks~\ref{rem:notgeneral} and Section~\ref{subsec:nonkilling}) between the ungaugeable Killing--commutant sector and the non-Killing sector in which massive gauge bosons and the Higgs do arise.
\end{enumerate}


%=====================================================================
\section{Conclusion}
\label{sec:conclusion}
%=====================================================================

We have proven that the $U(1)_7$ symmetry of the Dirac operator on Jensen-deformed $\mathrm{SU}(3)$, generated by the Kosmann derivative $K_7$, cannot be gauged within the framework of Connes' noncommutative geometry.  Two structurally distinct arguments establish this: the commutant obstruction (Proof~I, with its non-perturbative extension via the spectral theorem) and the categorical inner/outer classification (Proof~II, which is independent of the commutant property).  The Anderson--Higgs impossibility for $U(1)_7$ follows as a direct corollary.

The result rests on two complementary pillars.  The commutant property $[K_7, D_K] = 0$ simultaneously certifies $K_7$ as a symmetry generator (Kosmann--Lichnerowicz theorem) and obstructs its gauging (trivial one-form in $\Omega^1_D(\Alg)$) -- the isometry--obstruction duality of Proposition~\ref{prop:duality}.  Independently, the categorical position of $K_7$ as a diffeomorphism generator -- outer for the commutative algebra $C^\infty(\mathrm{SU}(3))$ -- excludes it from the gauge group $\Inn(\Alg_F)$ regardless of its commutant status.

The result is a structural wall in the space of mass-generation mechanisms for Kaluza--Klein spectral triples.  Any symmetry arising from an isometry of the internal manifold that commutes with the internal Dirac operator is permanently excluded from the gauge sector.  The spectral action, by its dependence on the spectrum of $D$ alone, inherits and enforces this exclusion automatically.  Crucially, the exclusion is sharp rather than general: it isolates the Killing--commutant sector, leaving the non-Killing sector -- where massive gauge bosons and the Higgs of the spectral action reside -- untouched (Sections~\ref{subsec:nonkilling} and \ref{subsec:saouter}).

Several directions for further investigation present themselves.  First, the classification of \emph{all} commutant elements of $D_K(\tau)$ -- not just $K_7$ -- would determine the full space of ungaugeable symmetries for each value of $\tau$.  Second, the fluctuated operator $D_{\mathrm{phys}}$ breaks $[iK_7, D_{\mathrm{phys}}] \neq 0$ (Remark~\ref{rem:dphys}); while this does not enable gauging (by Proof~II), it raises the question of whether the broken $U(1)_7$ has observable consequences in the low-energy effective theory.  Third, the interplay between ungaugeability and the persistence of the BCS Goldstone mode may constrain the physical spectrum of Kaluza--Klein compactifications on group manifolds.


%=====================================================================
% DATA AND CODE AVAILABILITY
%=====================================================================
\section*{Data and Code Availability}

The numerical verifications reported here -- the commutant identities of Proposition~\ref{prop:commutant} and Corollary~\ref{cor:allorders}, the charge decomposition, and the fluctuated-commutant value of Remark~\ref{rem:dphys} -- were produced by explicit matrix computations on the sixteen-dimensional singlet sector of the Peter--Weyl decomposition of $L^2(\mathrm{SU}(3),\Sigma)$, with the Dirac operator $D_K(\tau)$, the Kosmann operator $K_7$, and the finite-algebra generators constructed as described in Section~\ref{sec:prelim}.  The corresponding scripts and data are available from the authors on reasonable request.  All numerical statements in the paper are corroborations of analytically established results (Proposition~\ref{prop:commutant}, Corollary~\ref{cor:allorders}, Proofs~I and~II); the mathematical content does not depend on the computations.


%=====================================================================
% BIBLIOGRAPHY
%=====================================================================

\bibliographystyle{plainnat}
\bibliography{references}

\end{document}
